% ═══════════════════════════════════════════════════════════════
% LEHRERHINWEISE: Das Ohmsche Gesetz (Klasse 8)
% Kurzübersicht, Stundenverlauf, Differenzierung, Lösungen
% ═══════════════════════════════════════════════════════════════

\documentclass[11pt, a4paper]{article}

\usepackage[utf8]{inputenc}
\usepackage[T1]{fontenc}
\usepackage[ngerman]{babel}
\usepackage{helvet}
\renewcommand{\familydefault}{\sfdefault}

\usepackage[margin=2cm, top=2cm, bottom=2cm]{geometry}
\usepackage{amsmath, amssymb}
\usepackage{siunitx}
\sisetup{locale = DE, per-mode = symbol}

\usepackage{enumitem}
\usepackage{tabularx}
\usepackage{booktabs}
\usepackage{array}
\usepackage{longtable}

\usepackage{xcolor}
\usepackage{tcolorbox}
\usepackage{fancyhdr}
\usepackage{lastpage}

\setlist{nosep, leftmargin=1.5em}

% ═══════════════════════════════════════════════════════════════
% FARBEN
% ═══════════════════════════════════════════════════════════════

\definecolor{physik}{HTML}{2c5aa0}
\definecolor{grau}{HTML}{666666}
\definecolor{hellgrau}{HTML}{f0f0f0}
\definecolor{success}{HTML}{2a7a4b}
\definecolor{warning}{HTML}{b35c00}
\definecolor{error}{HTML}{c62828}

% ═══════════════════════════════════════════════════════════════
% KOPF-/FUSSZEILE
% ═══════════════════════════════════════════════════════════════

\pagestyle{fancy}
\fancyhf{}
\renewcommand{\headrulewidth}{0.4pt}
\lhead{\small\textcolor{physik}{\textbf{Physik Kl. 8}} -- Das Ohmsche Gesetz -- \textbf{Lehrerhinweise}}
\rhead{\small LH-01}
\cfoot{\small\thepage/\pageref{LastPage}}

% ═══════════════════════════════════════════════════════════════
% BOXEN
% ═══════════════════════════════════════════════════════════════

\tcbuselibrary{skins,breakable}

\newtcolorbox{infobox}[1][]{
    colback=hellgrau,
    colframe=physik,
    boxrule=1pt,
    arc=0pt,
    left=8pt, right=8pt, top=6pt, bottom=6pt,
    fonttitle=\bfseries,
    title=#1,
    breakable
}

\newtcolorbox{warnbox}[1][]{
    colback=white,
    colframe=warning,
    boxrule=1pt,
    arc=0pt,
    left=8pt, right=8pt, top=6pt, bottom=6pt,
    fonttitle=\bfseries,
    title=#1,
    breakable
}

\newtcolorbox{tippbox}{
    colback=white,
    colframe=success,
    boxrule=1pt,
    arc=0pt,
    left=8pt, right=8pt, top=6pt, bottom=6pt,
    fonttitle=\bfseries,
    title=Tipp,
    breakable
}

% ═══════════════════════════════════════════════════════════════
% DOKUMENT
% ═══════════════════════════════════════════════════════════════

\begin{document}

% ═══════════════════════════════════════════════════════════════
% TITEL
% ═══════════════════════════════════════════════════════════════

\begin{center}
{\Large\textbf{\textcolor{physik}{Lehrerhinweise: Das Ohmsche Gesetz}}}\\[0.3em]
{\normalsize Physik Klasse 8 -- Stunde 1 von 4 -- Reihe ``Elektrischer Widerstand''}
\end{center}

\vspace{1em}

% ═══════════════════════════════════════════════════════════════
% 1. KURZÜBERSICHT
% ═══════════════════════════════════════════════════════════════

\section*{1. Kurzübersicht}

\begin{infobox}
\begin{tabular}{ll}
\textbf{Thema:} & Das Ohmsche Gesetz (Einführung) \\
\textbf{Klasse:} & 8 (Gesamtschule, Profil A: heterogen) \\
\textbf{Zeit:} & 45 Minuten \\
\textbf{Lernziele:} & Proportionalität $U \sim I$, Formel $U = R \cdot I$ anwenden \\
\textbf{Methoden:} & Demonstrationsexperiment, Schülerexperiment, differenziertes AB \\
\textbf{Material:} & Netzgerät, Widerstände (\SI{20}{\ohm}), Multimeter, AB-01-ohm \\
\end{tabular}
\end{infobox}

% ═══════════════════════════════════════════════════════════════
% 2. STUNDENVERLAUF
% ═══════════════════════════════════════════════════════════════

\section*{2. Stundenverlauf (5-Minuten-Raster)}

\begin{longtable}{|c|p{3cm}|p{7cm}|p{3cm}|}
\hline
\textbf{Min.} & \textbf{Phase} & \textbf{Inhalt} & \textbf{Material} \\
\hline
\endhead

0--5 & Einstieg & Demo: Lampe an Netzgerät, $U$ variieren. Impuls: ``Was passiert mit der Helligkeit?'' & Netzgerät, Lampe \\
\hline

5--10 & Überleitung & Schaltbild erklären, Experiment einführen, Gruppen einteilen & Tafel \\
\hline

10--15 & Experiment & Gruppen bauen auf, messen $I$ bei $U = \SI{2}{\volt}$ & Experimentiermaterial \\
\hline

15--20 & Experiment & Messungen bei $U = 4, 6, 8, \SI{10}{\volt}$ & Experimentiermaterial \\
\hline

20--25 & Experiment & Letzte Messungen, aufräumen & Experimentiermaterial \\
\hline

25--30 & Sicherung I & Werte sammeln, gemeinsam Diagramm erstellen, Proportionalität erkennen & Tafel, Doku-Kamera \\
\hline

30--35 & Übung & AB verteilen, Niveau erklären, Bearbeitung beginnen & AB-01-ohm \\
\hline

35--40 & Übung & Weiterarbeit am AB, L geht herum und hilft & AB-01-ohm \\
\hline

40--45 & Sicherung II & Musterlösungen besprechen, Hefteintrag: Formel + Merksatz & Tafel, Heft \\
\hline
\end{longtable}

% ═══════════════════════════════════════════════════════════════
% 3. DIFFERENZIERUNG
% ═══════════════════════════════════════════════════════════════

\section*{3. Differenzierung}

\begin{tabular}{|l|p{4cm}|p{4cm}|p{4cm}|}
\hline
& \textbf{Basis (*)} & \textbf{Standard (**)} & \textbf{Erweitert (***)} \\
\hline
\textbf{Ziel} & Qualitatives Verständnis & Formel anwenden & Formel umstellen, Transfer \\
\hline
\textbf{Aufgaben} & 1--2 (5 BE) & 1--4 (14 BE) & Alle (25 BE) \\
\hline
\textbf{Hilfen} & Formelkarte, Beispiel auf AB & Hinweise bei Umstellung & Nur Impulse \\
\hline
\textbf{Experiment} & Messen, Werte eintragen & + Diagramm interpretieren & + Proportionalität begründen \\
\hline
\textbf{Inklusion} & LRS: Zeitzugabe, größere Schrift & & Expertenrolle für Starke \\
\hline
\end{tabular}

\vspace{1em}

\begin{tippbox}
Bei heterogenen Gruppen: Leistungsstarke SuS als ``Messleiter'' einsetzen, die anderen notieren Werte.
\end{tippbox}

% ═══════════════════════════════════════════════════════════════
% 4. HÄUFIGE SCHÜLERFEHLER
% ═══════════════════════════════════════════════════════════════

\section*{4. Häufige Schülerfehler und Reaktionen}

\begin{warnbox}[Typische Fehler]
\begin{tabular}{|p{5cm}|p{8cm}|}
\hline
\textbf{Fehler} & \textbf{Reaktion / Hilfe} \\
\hline
\textbf{mA nicht umgerechnet} (200 mA statt 0,2 A) & ``Schau auf die Einheit -- wie viel ist 1 A in mA?'' \\
\hline
\textbf{Formel falsch umgestellt} ($I = U \cdot R$) & URI-Dreieck zeigen; ``Was steht oben allein?'' \\
\hline
\textbf{Einheiten vergessen} & ``Jede physikalische Größe braucht eine Einheit!'' \\
\hline
\textbf{Verwechslung $U$ und $I$} & ``Was ist der Antrieb? Was fließt?'' \\
\hline
\textbf{``Strom wird verbraucht''} & ``Der Strom fließt durch -- er kommt genauso viel raus wie rein.'' \\
\hline
\textbf{``$R$ ändert sich bei höherem $U$''} & ``$R$ ist eine Eigenschaft des Drahtes -- wie seine Länge.'' \\
\hline
\end{tabular}
\end{warnbox}

% ═══════════════════════════════════════════════════════════════
% 5. MATERIAL-CHECKLISTE
% ═══════════════════════════════════════════════════════════════

\section*{5. Material-Checkliste}

\begin{itemize}
    \item[$\Box$] Netzgerät (0--12 V, regelbar) für Demo
    \item[$\Box$] Glühlampe (12 V) für Demo
    \item[$\Box$] 6 Experimentiersätze (je Gruppe):
    \begin{itemize}
        \item Netzgerät
        \item Widerstand \SI{20}{\ohm} (fest)
        \item 2 Multimeter (für $U$ und $I$)
        \item Kabel
    \end{itemize}
    \item[$\Box$] AB-01-ohm.pdf kopiert (1x pro SuS, ggf. größer für LRS)
    \item[$\Box$] Formelkarte für Niveau (*)
    \item[$\Box$] Tafelanschrieb vorbereitet: Schaltbild, Tabelle, Diagramm-Achsen
\end{itemize}

% ═══════════════════════════════════════════════════════════════
% 6. LÖSUNGEN (KURZFORM)
% ═══════════════════════════════════════════════════════════════

\section*{6. Lösungen (Kurzform)}

\textit{Vollständige Lösungen: siehe ML-01-ohm.pdf}

\vspace{0.5em}

\begin{tabular}{|l|l|}
\hline
\textbf{Aufgabe} & \textbf{Lösung} \\
\hline
Formelbox & Volt (V), Widerstand, $\Omega$, Ampere (A) \\
\hline
Merksatz & höher, höher, konstant \\
\hline
Messwerte & 0,1 / 0,2 / 0,3 / 0,4 / 0,5 A \\
\hline
1 & verdoppelt, Stromstärke, proportional \\
\hline
2 & $U = \SI{10}{\ohm} \cdot \SI{2}{\ampere} = \SI{20}{\volt}$ \\
\hline
3a & $U = \SI{10}{\volt}$ \\
\hline
3b & $I = \SI{3}{\ampere}$ \\
\hline
3c & $R = \SI{20}{\ohm}$ \\
\hline
4 & $U \to V$, $I \to A$, $R \to \Omega$ \\
\hline
5 & mA-Umrechnung vergessen; richtig: \SI{10}{\volt} \\
\hline
6a & $I = \SI{0,3}{\ampere}$ \\
\hline
6b & Verdopplung (Proportionalität) \\
\hline
7 & $R$ ist Materialeigenschaft \\
\hline
\end{tabular}

% ═══════════════════════════════════════════════════════════════
% 7. TAFELANSCHRIEB
% ═══════════════════════════════════════════════════════════════

\section*{7. Tafelanschrieb (Vorschlag)}

\begin{infobox}[Tafel links: Schaltbild]
\begin{center}
\begin{verbatim}
    +---[R=20Ω]---+
    |             |
   (U)           (A)
    |             |
    +-----[~]-----+
      Netzgerät
\end{verbatim}
\end{center}
\end{infobox}

\begin{infobox}[Tafel Mitte: Messwerte + Diagramm]
\begin{tabular}{|c|c|c|c|c|c|}
\hline
$U$ in V & 2 & 4 & 6 & 8 & 10 \\
\hline
$I$ in A & 0,1 & 0,2 & 0,3 & 0,4 & 0,5 \\
\hline
\end{tabular}

\vspace{0.5em}
$\Rightarrow$ Gerade durch den Ursprung = Proportionalität
\end{infobox}

\begin{infobox}[Tafel rechts: Formel + Merksatz]
\textbf{Das Ohmsche Gesetz:}

\vspace{0.5em}
\begin{center}
{\Large $U = R \cdot I$}
\end{center}

\vspace{0.5em}
Je höher die Spannung, desto höher die Stromstärke.

Der Widerstand $R$ bleibt konstant.
\end{infobox}

% ═══════════════════════════════════════════════════════════════
% 8. AUSBLICK
% ═══════════════════════════════════════════════════════════════

\section*{8. Ausblick: Folgestunden}

\begin{tabular}{|c|l|l|}
\hline
\textbf{Std.} & \textbf{Thema} & \textbf{Schwerpunkt} \\
\hline
2 & Widerstand messen & Kennlinien aufnehmen, ohmsch vs. nicht-ohmsch \\
\hline
3 & Widerstände schalten & Reihen- und Parallelschaltung, $R_{\text{ges}}$ \\
\hline
4 & Anwendungen & Spannungsteiler, Vorwiderstand LED \\
\hline
\end{tabular}

\vspace{1em}

\begin{tippbox}
In Stunde 2: Glühlampe als nicht-ohmscher Widerstand (Kennlinie krümmt sich).
SuS entdecken, dass $R$ bei Glühlampe \textit{nicht} konstant ist.
\end{tippbox}

% ═══════════════════════════════════════════════════════════════
% 9. NOTIZEN
% ═══════════════════════════════════════════════════════════════

\section*{9. Notizen (nach der Durchführung)}

\begin{tabular}{|p{7cm}|p{7cm}|}
\hline
\textbf{Was hat gut funktioniert?} & \textbf{Was muss angepasst werden?} \\
\hline
& \\[3cm]
\hline
\end{tabular}

\vspace{1em}

\textbf{Zeiteinhaltung:} $\Box$ okay \quad $\Box$ zu knapp \quad $\Box$ Zeit übrig

\textbf{Differenzierung:} $\Box$ okay \quad $\Box$ anpassen (wie: \rule{5cm}{0.4pt})

% ═══════════════════════════════════════════════════════════════
% ENDE
% ═══════════════════════════════════════════════════════════════

\end{document}
